\chapter[Stato incrementale]{Stato incrementale e lavoro\\riutilizzabile}

\begin{flushleft}
\textbf{Origine:} Week 002, esercizi 3 e 4\\
\textbf{Contesto:} batch contigui entro un budget e registro dinamico di modelli\\
\textbf{Esito della review:} output corretti, vincoli di scala non rispettati\\
\textbf{Pattern:} sliding window con somma corrente; query online incrementali
\end{flushleft}

\section{Principio comune: conservare lo stato}

I due problemi della settimana hanno formulazioni differenti, ma condividono
la stessa struttura algoritmica. Una sequenza di operazioni consecutive
modifica soltanto una piccola parte dello stato precedente:

\begin{itemize}
  \item nella finestra sul budget entra un costo dall'estremo destro e possono
        uscire alcuni costi dall'estremo sinistro;
  \item nel registro cambia la frequenza di un solo modello per operazione.
\end{itemize}

Ricalcolare ogni volta l'intera finestra o l'intero registro produce risultati
corretti, ma perde l'informazione già disponibile. L'obiettivo dell'approccio
incrementale è mantenere un riassunto sufficiente dello stato e aggiornarlo
soltanto in corrispondenza della modifica corrente.

\begin{principlebox}{Domanda diagnostica}
Se due iterazioni consecutive condividono quasi tutti i dati, quale quantità
può essere aggiornata aggiungendo ciò che entra e rimuovendo ciò che esce?
Questa domanda separa spesso una soluzione quadratica da una lineare.
\end{principlebox}

\section{Sliding window con vincolo sulla somma}

\subsection{Modello e precondizioni}

Sia $C=(c_0,\ldots,c_{n-1})$ una sequenza di costi non negativi e sia
$B\geq 0$ il budget. Si cerca la massima lunghezza di un segmento contiguo
$C[l:r]$ tale che
\[
  \sum_{i=l}^{r} c_i \leq B.
\]

La non negatività è una precondizione essenziale. Se la somma supera il
budget, estendere ulteriormente la finestra non può ridurla; per renderla di
nuovo valida occorre avanzare l'estremo sinistro. Con valori negativi questa
monotonia non sarebbe garantita e servirebbe un algoritmo diverso.

\subsection{Submission originale e diagnosi}

La submission seguente è funzionalmente corretta sugli input ammessi. La
variabile \texttt{i} seleziona ogni possibile inizio e \texttt{j} estende il
segmento finché il budget è rispettato. Il problema è che \texttt{j} e
\texttt{tot\_budget} vengono reinizializzati per ogni nuovo \texttt{i}; i
costi già sommati vengono quindi visitati molte volte.

Il codice conserva integralmente la submission; nella sola stringa d'errore,
\texttt{può} è reso come \texttt{puo'} per compatibilità tipografica con il
motore dei riquadri. La modifica non ha alcun effetto sulla logica.

\begin{inefficientcode}{Submission originale - output corretti, costo quadratico}
def longest_batch_within_budget(costs, budget):
    """Restituisce la massima lunghezza contigua entro il budget."""
    len_series=0
    if budget <0:
        raise ValueError("Il budget non puo' essere negativo")

    for i in range(len(costs)):
        j=0
        left=costs[i]
        tot_budget=left
        n_serie=0
        while tot_budget<=budget:
            n_serie+=1
            j+=1
            if i+j<len(costs):
                tot_budget+=costs[i+j]
            else:
                break
        len_series =max(len_series, n_serie)
    return len_series
\end{inefficientcode}

Nel caso peggiore, ad esempio con tutti i costi uguali a zero, il numero di
visite è
\[
  n+(n-1)+\cdots+1=\frac{n(n+1)}{2}=\Theta(n^2).
\]
L'espressione \texttt{i+j} calcola un indice in tempo costante, ma non riduce
il numero complessivo di elementi riesaminati. Nei benchmark della review, il
passaggio da 4.000 a 8.000 elementi ha moltiplicato il tempo per circa quattro,
coerentemente con la crescita quadratica.

\clearpage
\thispagestyle{fancy}
\noindent\rule{0pt}{1cm}\par
\subsection{Algoritmo lineare e invariante}

La finestra efficiente mantiene gli indici \texttt{left} e \texttt{right} e
una somma corrente. Quando \texttt{right} avanza, il nuovo costo viene
aggiunto. Se la somma supera il budget, si sottraggono in ordine i costi che
escono da sinistra.

\begin{principlebox}{Invariante della somma corrente}
Al termine del ciclo \texttt{while}, vale
\[
\texttt{current\_sum}=\sum_{i=\texttt{left}}^{\texttt{right}}c_i\leq B,
\]
e \texttt{left} è il più piccolo indice che rende valida la finestra con
estremo destro \texttt{right}.
\end{principlebox}

\begin{efficientcode}{Implementazione di riferimento - sliding window O(n)}
def longest_batch_within_budget(costs, budget):
    """Restituisce la massima lunghezza contigua entro il budget."""
    if budget < 0:
        raise ValueError("Il budget non puo' essere negativo")

    left = 0
    current_sum = 0
    best = 0

    for right, cost in enumerate(costs):
        current_sum += cost

        while current_sum > budget:
            current_sum -= costs[left]
            left += 1

        best = max(best, right - left + 1)

    return best
\end{efficientcode}

\subsection{Argomento di correttezza}

Prima di elaborare un elemento la somma descrive esattamente la finestra
precedente. L'aggiunta di $c_{\texttt{right}}$ estende tale finestra di un
elemento. Se il budget è superato, ogni iterazione del \texttt{while} sottrae
esattamente il costo escluso e avanza \texttt{left}; l'uguaglianza tra somma
memorizzata e somma del segmento è quindi preservata.

Quando il ciclo termina, la finestra è valida. Inoltre \texttt{left} è minimo:
il costo precedente è stato rimosso soltanto mentre la somma superava $B$.
La lunghezza \texttt{right-left+1} è pertanto la massima lunghezza valida con
quell'estremo destro. Considerando tutti gli estremi destri e conservando il
massimo si ottiene la soluzione globale.

Ogni elemento entra una volta tramite \texttt{right} ed esce al massimo una
volta tramite \texttt{left}. Il tempo è $O(n)$ e lo spazio ausiliario è
$O(1)$.

\section{Query online sul numero di modelli attivi}

\subsection{Modello del registro}

Sia $f(m)$ il numero corrente di deployment del modello $m$. Le operazioni
possono incrementare o decrementare una frequenza, richiederne il valore o
chiedere il numero di modelli distinti attivi:
\[
  A=\left|\{m\mid f(m)>0\}\right|.
\]

Una query è detta \emph{online} quando deve essere risolta nello stato prodotto
da tutte le operazioni precedenti, senza conoscere o riorganizzare quelle
successive.

\subsection{Submission originale e costo delle query}

La submission gestisce correttamente registrazioni, ritiri, conteggi e
transizioni a zero. La query \texttt{distinct}, tuttavia, visita ogni chiave
del dizionario per ricostruire $A$.

\begin{inefficientcode}{Submission originale - distinct ricalcolato integralmente}
def process_model_registry(initial_models, operations):
    """Elabora aggiornamenti e query sul registro dei modelli."""
    result=[]
    registry_model={}
    for model in initial_models:
        if model not in registry_model:
            registry_model[model]=0
        registry_model[model]+=1

    for command in operations:
        if command[0]=="count":
            model=command[1]
            count=registry_model.get(model,0)
            result.append(count)
            continue
        if command[0]=="distinct":
            count=0
            for model in registry_model:
                if registry_model.get(model)!=0:
                    count+=1
            result.append(count)
            continue
        if command[0]=="retire":
            model=command[1]
            if model in registry_model and registry_model.get(model)>0:
                registry_model[model]-=1
            continue
        if command[0]=="register":
            model=command[1]
            registry_model[model]=registry_model.get(model,0)+1
            continue
    return result
\end{inefficientcode}

Se il registro contiene $m$ chiavi e si eseguono $q$ query
\texttt{distinct}, il costo può raggiungere $O(qm)$. Con $m$ e $q$ entrambi
proporzionali alla dimensione dell'input, il comportamento è quadratico.

\subsection{Aggregato incrementale e transizioni di soglia}

Una singola operazione modifica la frequenza di un solo modello. Il valore $A$
non cambia per ogni incremento o decremento, ma soltanto quando una frequenza
attraversa la soglia zero:

\begin{itemize}
  \item $0\rightarrow1$: un modello diventa attivo, quindi $A$ aumenta;
  \item $1\rightarrow0$: un modello cessa di essere attivo, quindi $A$
        diminuisce;
  \item $k\rightarrow k+1$ con $k>0$ e $k\rightarrow k-1$ con $k>1$: $A$
        resta invariato.
\end{itemize}

\begin{principlebox}{Invariante del registro}
Dopo ogni operazione, \texttt{active\_models} è uguale al numero di chiavi del
dizionario la cui frequenza è strettamente positiva. La query
\texttt{distinct} può quindi rispondere in $O(1)$.
\end{principlebox}

\begin{efficientcode}{Implementazione di riferimento - query tutte O(1) in media}
def process_model_registry(initial_models, operations):
    """Elabora aggiornamenti e query sul registro dei modelli."""
    frequencies = {}
    active_models = 0
    results = []

    for model in initial_models:
        current = frequencies.get(model, 0)
        if current == 0:
            active_models += 1
        frequencies[model] = current + 1

    for operation in operations:
        command = operation[0]

        if command == "count":
            results.append(frequencies.get(operation[1], 0))

        elif command == "distinct":
            results.append(active_models)

        elif command == "register":
            model = operation[1]
            current = frequencies.get(model, 0)
            if current == 0:
                active_models += 1
            frequencies[model] = current + 1

        elif command == "retire":
            model = operation[1]
            current = frequencies.get(model, 0)
            if current > 0:
                frequencies[model] = current - 1
                if current == 1:
                    active_models -= 1

    return results
\end{efficientcode}

\subsection{Correttezza e complessità}

Durante l'inizializzazione, \texttt{active\_models} viene incrementato
esattamente alla prima occorrenza di ogni modello, quindi l'invariante vale
prima delle query. Un'operazione \texttt{count} o \texttt{distinct} non cambia
lo stato. Una registrazione modifica $A$ soltanto se il conteggio precedente è
zero; un ritiro valido lo modifica soltanto se il conteggio precedente è uno.
In ogni caso l'invariante è preservato per induzione.

Assumendo accessi hash $O(1)$ in media, inizializzazione e operazioni richiedono
tempo $O(n+q)$, dove $n$ è il numero di deployment iniziali e $q$ il numero di
operazioni. Lo spazio è $O(m+r)$: $m$ modelli distinti memorizzati e $r$
risposte prodotte.

\section{Segnali di riconoscimento ed edge case}

\begin{itemize}
  \item \textbf{Finestra contigua e valori non negativi:} una soglia sulla
        somma suggerisce una finestra espandibile e restringibile.
  \item \textbf{Query ripetute su un aggregato:} se ogni update tocca un solo
        elemento, verificare se l'aggregato può essere aggiornato localmente.
  \item \textbf{Soglie di stato:} conteggi di elementi attivi, non vuoti o
        disponibili cambiano spesso soltanto attraversando zero.
  \item \textbf{Sequenza vuota:} la finestra restituisce zero; il registro
        parte con zero modelli attivi.
  \item \textbf{Ritiro di un elemento assente:} non cambia frequenze né
        aggregato.
  \item \textbf{Costi uguali a zero:} possono restare tutti nella finestra e
        devono contribuire alla lunghezza.
\end{itemize}

\begin{definitionbox}{Stato di apprendimento}
I due invarianti sono stati analizzati e compresi durante la review, ma non
sono ancora masterizzati. La mastery richiede implementazioni autonome,
lineari e completate entro il timebox su formulazioni differenti.
\end{definitionbox}
